\subsection{System Overview}

The proposed pipeline operates in two sequential stages at inference time:

\begin{enumerate}
    \item \textbf{Hand Detection (MediaPipe):} Each incoming video frame is passed to a MediaPipe model that produces bounding boxes around detected hand regions.
    \item \textbf{Letter Classification (ResNet-50):} The detected hand crop is extracted, resized, and passed to the classification model, which returns a probability distribution over 26 ASL letter classes.
\end{enumerate}

The predicted letter and bounding box are overlaid on the original frame for visual feedback.

\subsection{Hand Detection Module}

The \texttt{HandDetector} class wraps the MediaPipe HandLandmarker Tasks API,
operating in \texttt{VIDEO} mode to enable frame-to-frame tracking via MediaPipe's
internal Kalman filter. This mode requires strictly monotonically increasing
real wall-clock millisecond timestamps per frame. A common implementation error
is to pass bare integer counters (1, 2, 3, \ldots); doing so mis-models velocity
in the Kalman filter and worsens landmark jitter. Timestamps are derived here
from \texttt{time.monotonic\_ns()} to satisfy this requirement correctly.

Detection thresholds are set above MediaPipe's defaults: minimum hand detection
and presence confidence at 0.65 (vs.\ default 0.50), and tracking confidence at
0.55. The looser defaults are too permissive for static-sign ASL, where ambiguous
partially-occluded hands during transitions cause false detections.

Per-frame, the detector returns a list of up to two hand results. For each hand,
the raw landmark bounding box is computed from the min/max of all 21 normalized
landmark coordinates, then expanded by 35\% in each dimension (\texttt{expansion=0.35})
to provide context around the hand edges. The expanded box is then smoothed
across frames using exponential smoothing weighted at $0.3 \cdot \text{prev} +
0.7 \cdot \text{new}$, keeping the bbox responsive to motion. Crucially, motion
magnitude is computed on the \emph{raw} (pre-smoothing) center against the
previous smoothed center. Computing motion after smoothing — as a naive
implementation would do — causes a fast-moving hand to appear stable because
the smoothed center barely shifts, defeating the stability gate entirely.
Motion is measured as Euclidean (L2) distance; Manhattan (L1) distance
under-reports diagonal motion by approximately 30\%.

Per-hand tracking state (previous bbox and center) is keyed by hand index, not
stored as a single shared variable. A single shared variable fails silently the
moment two hands appear in frame. When a hand disappears, its tracking state is
evicted immediately so the next detection initializes fresh rather than snapping
to stale coordinates. If the clamped ROI crop has zero area after boundary
clamping, that detection is skipped rather than passing an empty array to the
classifier.

At inference time, when multiple hands are detected, the detection with the
highest handedness classifier score is selected for classification:

\begin{equation}
    \text{det}^* = \arg\max_{d \in \mathcal{D}}\; d.\text{score}
\end{equation}

The \texttt{HandDetection} dataclass returned per hand carries: the smoothed
pixel bbox $(x_1, y_1, x_2, y_2)$, handedness label and score, the 21 raw
normalized landmarks, the BGR ROI crop, a boolean stability flag, and the raw
motion magnitude in pixels.

\subsection{Classification Model (SLD)}

The Sign Language Detector (\textbf{SLD}) model consists of a ResNet-50 backbone
with all parameters frozen, followed by a trainable lightweight classification head.
The backbone's final average pooling layer produces a 2048-dimensional feature
vector. The original fully-connected layer is replaced entirely. The classification
head is defined as:

\begin{equation}
    \hat{y} = \text{head}(\text{backbone}(x))
\end{equation}

\begin{equation}
    \text{head}(z) = W_2 \cdot \text{GELU}(\text{Dropout}(\text{LayerNorm}(z) \cdot W_1))
\end{equation}

with $W_1 \in \mathbb{R}^{2048 \times 512}$ and $W_2 \in \mathbb{R}^{512 \times 29}$,
supporting all 29 output classes (26 letters plus \texttt{del}, \texttt{nothing},
\texttt{space}).

\textbf{LayerNorm} is used at the backbone-to-head boundary rather than BatchNorm.
BatchNorm at this position is sensitive to batch size and behaves differently at
train vs.\ eval time, which can destabilize fine-tuning when the backbone is frozen
and only the head is updated. LayerNorm normalizes per-sample and is invariant to
batch size. \textbf{GELU} activation is chosen for its smoothness and non-zero
gradient at negative values compared to ReLU. \textbf{Dropout} at $p = 0.3$
regularizes the intermediate 512-dimensional representation.

At inference, the model is loaded from a checkpoint via \texttt{torch.load} with
\texttt{weights\_only=False}, explicitly stated to suppress the PyTorch 2.x
deprecation warning and to document intent — the checkpoint stores optimizer and
scheduler state in addition to model weights and cannot be loaded in
weights-only mode.

\subsection{Training Configuration}

Training is governed by a dataclass-based configuration (\texttt{TRAINING\_CONFIG})
that encapsulates all hyperparameters for reproducibility and supports runtime
overrides via command-line arguments. Key settings are summarized in
Table~\ref{tab:config}.

The dataset is split using \texttt{random\_split} into train (70\%), validation
(20\%), and test (10\%) subsets. A custom \texttt{\_TransformSubset} wrapper is
applied to each split so that training and evaluation transforms can differ
without modifying the underlying dataset object. A naive approach of overwriting
\texttt{subset.dataset} causes the Subset's indices to point into a
wrongly-sized dataset and raises an \texttt{IndexError}; wrapping the Subset
itself avoids this. Training images receive data augmentation; validation and
test images receive only deterministic resize and normalization.

Per-batch metrics — training loss, accuracy, and gradient norm — are logged to
W\&B at every step via \texttt{wandb.log(..., step=self.global\_step)}, giving
fine-grained visibility into optimization dynamics beyond epoch-level summaries.
At the epoch level, mean gradient norm and gradient clipping ratio (fraction of
batches where the norm exceeded the clip threshold) are also reported, providing
a direct signal on how aggressively the optimizer is being constrained.

\begin{table}[h]
\caption{Training Hyperparameters}
\label{tab:config}
\centering
\begin{tabular}{ll}
\toprule
\textbf{Parameter} & \textbf{Value} \\
\midrule
Optimizer & SGD (Nesterov) \\
Learning Rate & $1 \times 10^{-2}$ \\
Momentum & 0.9 \\
Weight Decay & $5 \times 10^{-4}$ \\
Scheduler & Cosine Annealing + Warm-Up \\
Warm-Up Epochs & 10 \\
$T_{\max}$ & matched to epoch count \\
$\eta_{\min}$ & $1 \times 10^{-4}$ \\
Max Epochs & 200 (3 in practice) \\
Batch Size & 128 \\
Gradient Clip Norm & 1.0 \\
Early Stopping Patience & 20 \\
\bottomrule
\end{tabular}
\end{table}

Due to Google Colab compute constraints and session limitations, training was 
conducted for only 3 epochs rather than the full 200 configured. The \texttt{T\_max} 
and epoch count parameters were overridden via command-line arguments at runtime. 
Despite this, as shown in Section~\ref{sec:results}, the model converged to 
competitive accuracy within these three epochs, suggesting the dataset and 
backbone are well-matched for this task.

\subsubsection{Learning Rate Schedule}

A cosine annealing schedule with linear warm-up is used. During the warm-up phase (epochs $0$ to $T_w - 1$), the learning rate increases linearly:

\begin{equation}
    \eta_t = \eta_{\text{base}} \cdot \frac{t+1}{T_w}
\end{equation}

After warm-up, cosine decay is applied:

\begin{equation}
    \eta_t = \eta_{\min} + \frac{\eta_{\text{base}} - \eta_{\min}}{2} \left(1 + \cos\left(\pi \cdot \frac{t - T_w}{T_{\max} - T_w}\right)\right)
\end{equation}

This prevents early training instability while allowing the learning rate to settle to a low value near convergence.

\subsubsection{Gradient Clipping}

Gradient norm clipping at threshold $\tau = 1.0$ is applied after each backward pass to prevent exploding gradients:

\begin{equation}
    \mathbf{g} \leftarrow \mathbf{g} \cdot \min\!\left(1,\ \frac{\tau}{\|\mathbf{g}\|_2}\right)
\end{equation}

\subsubsection{Early Stopping}

Training is halted if validation accuracy does not improve for 20 consecutive
epochs. The best-performing checkpoint — saving model, optimizer, scheduler
state, and best accuracy — is written to disk whenever a new peak is reached.
At the end of training, this best checkpoint is reloaded before final test
evaluation, ensuring test metrics reflect the best generalization point rather
than the last epoch.

\subsection{Experiment Tracking}

All training runs are logged via Weights \& Biases (W\&B) \cite{wandb}. Per-batch
metrics include training loss, training accuracy, and gradient norm. Per-epoch
metrics include validation loss and accuracy, learning rate, mean gradient norm,
and gradient clipping ratio. Final test loss and accuracy are logged as summary
metrics at run completion. The full \texttt{TRAINING\_CONFIG} — including
optimizer kwargs, scheduler kwargs, and dataloader configs — is serialized into
the W\&B run config for exact reproducibility.
